\clearpage\begingroup\SourceLeaf{625b}
\ArithmeticHead{70}{L}{PHILIPP MAENNCHEN, GAUSS ALS ZAHLENRECHNER.}
\renewcommand{\thefootnote}{\arabic{footnote})}\setcounter{footnote}{0}
\begin{center}\textls[120]{Anhang.}\end{center}
\begin{center}[1.]\\[3pt]
\textsc{Bessel} an \textsc{Gauss}. Bremen, 8. Sept. 1805.\end{center}
\noindent\hspace*{\parindent}. . . »Ich habe schon einen kleinen Anfang bei den Rechnungen gemacht.
Die Logarithmen der Brüche, die Sie mir mitteilen, sind einige Mal fehlerhaft:
\[
\begin{array}{rclcl}
\log\dfrac{1}{528}&\text{muss sein}&7{,}2773661&\text{statt}&7{,}3151546\\[5pt]
\log\dfrac{121}{120}&\text{muss sein}&0{,}0036041&\text{statt}&0{,}0036042\\[5pt]
\log\dfrac{49}{48}&\text{muss sein}&0{,}0089548&\text{statt}&0{,}0089549\\[5pt]
\log\dfrac{9}{8}&\text{muss sein}&0{,}0511525&\text{statt}&0{,}0511225.
\end{array}
\]
\clearpage\SourceLeaf{626}\ArithmeticHead{71}{R}{ANHANG.}
\noindent\hspace*{\parindent}Der letzte Logarithmus ist nur verschrieben; das mir mitgeteilte Exempel
enthält ihn schon richtig. Ich zeige Ihnen diese Differenzen an, weil sie auf
die Rechnungen, die Sie sich selbst vorbehalten, vielleicht Einfluss haben
könnten; augenscheinlich haben Sie statt $\log 528$ $\log 484$ angenommen.« . . .
\begin{center}[2.]\\[3pt]
\textsc{Bessel} an \textsc{Gauss}. Bremen, 19. September 1805.\end{center}
\noindent\hspace*{\parindent}Einliegend erhalten Sie, mein hochgeschätztester Freund, eine Tafel, die
die Koeffizienten enthält, deren Berechnung Sie mir übertrugen. Ich füge
nur die Bemerkung hinzu, dass die zweite Columne ganz nach Ihrer Rechnung
beibehalten wurde, und dass die rot geschriebenen Zahlen nach der meinigen
angesetzt sind; nur einmal (bei $30^0\,12'$) nahm ich das Mittel aus Ihrer und
meiner Rechnung, denn nur da zeigte sich ein Unterschied in der letzten
Dezimalstelle von 2 oder fast $2\frac{1}{2}$; wenn es Ihnen auf eine oder zwei Einheiten
der letzten Ziffer ankommt, so schlage ich Ihnen die Revision Ihrer
Rechnung für diesen Wert von $W$ vor. Ich halte übrigens die Koeffizienten
für richtig, welches jedoch nur in so fern wahr sein kann, in so fern man
die Anhäufung kleiner Fehler in den Logarithmen von $\frac{1}{1+f^2}$ oder $2\log\cos W$,
welche sich auf die höheren Koeffizienten äussern muss, vernachlässigt. Wenn
Sie, verehrungswürdiger Freund, nun einmal wieder eine Rechnung auszuführen
haben, der Sie gern überhoben wären, so sprechen Sie nur immer
damit bei mir vor; für die kleine Mühe finde ich reichlichen Ersatz in dem
Gedanken, für Sie zu arbeiten, und dann auch in dem Briefe, den Sie mir
dabei schreiben werden, wäre er auch nur halb so lehrreich als der letzte. . . .

[Eine Anlage enthält die Tafel der Koeffizienten zur Entwickelung von
$(a^2+a'^2-2aa'\cos\varphi)^{-\frac{1}{2}}$ in eine unendliche Reihe.]
\begin{center}[3.]\\[3pt]
\textsc{Gauss} an \textsc{Bessel}. Braunschweig, den 7. Oktober 1805.\end{center}
\noindent\hspace*{\parindent}Ich bin Ihnen, mein teurer Freund, noch den allerverbindlichsten Dank
schuldig für die gütige Übernehmung der beschwerlichen Rechnung. Ich würde
\clearpage\SourceLeaf{627}\ArithmeticHead{72}{L}{PHILIPP MAENNCHEN, GAUSS ALS ZAHLENRECHNER.}
\setcounter{footnote}{0}
\noindent Ihnen diesen schon weit früher dargebracht haben, wenn nicht eine jetzt vorgehende
Veränderung meiner häuslichen Lage\begingroup[\footnotemark]\renewcommand{\thefootnote}{[\arabic{footnote})}\footnotetext{Am 9. Oktober 1805 fand \textsc{Gauss}’ Vermählung mit \textsc{Johanna Osthoff} statt.]}\endgroup\ mich in so mancherlei Zerstreuungen
verwickelt hätte, dass ich an das Ordnen verschiedener Bemerkungen,
womit ich meinen Dank gern begleiten wollte, nicht habe denken können.
Auch jetzt kann ich dies noch nicht; um aber doch nicht gar zu lange zu
schweigen, will ich mit Vorbehalt einer künftigen ausführlicheren Unterhaltung,
namentlich über die Art, wie ich die beiden ersten Koeffizienten berechne,
nur ein Paar Anmerkungen hinzusetzen.

Zuerst danke ich Ihnen verbindlichst dafür, dass Sie mich auf die fehlerhaften
Logarithmen aufmerksam gemacht haben. Bei $\log\frac{1}{528}$ war wirklich
ein Versehen begangen, das wahrscheinlich mir einige unnütze Arbeit verursacht
hätte, wenn ich nicht glücklicherweise bei Empfang Ihrer Warnung
mit jener Rechnung noch keinen Anfang gemacht hätte. Bei $\frac{121}{120}$ und $\frac{49}{48}$
wird aber der Unterschied wohl nur daher rühren, dass ich diese Logarithmen
aus grösseren Tafeln genommen habe (gewiss kann ich’s nicht sagen, weil in
diesem Augenblicke der \textit{Thesaurus} von \textsc{Vega} noch eingepackt liegt), dieser
Unterschied ist übrigens von gar keiner Bedeutung. Ebenso würde auch der
Unterschied bei $30^0\,12'$, wo ich in meiner wieder nachgesehenen Rechnung
keinen Fehler finde, gleich um eine Einheit kleiner werden, wenn der doppelte
Logarithmus des Cosinus aus grösseren Tafeln abgeleitet wäre.

Ich bin übrigens in meinem Anteile an der Arbeit nun auch so weit
gekommen, dass ich sowohl die Rechnung für $\log B'',\,B'''$ etc. für die übrigen
Werte von $W$ ganz vollendet habe, als auch schon das ganze Interpolationsgeschäft
für $\log B^0,\,B',\,B''$ und $B'''$. Bei dieser freilich auch beschwerlichen
Arbeit kann ich indess um so weniger von fremder Hülfe Gebrauch machen,
da nicht viel mehr dabei zu tun ist, als das blosse Schreiben. Ich bediene
mich dabei der Formel
\[
\varphi(x)=\frac{9}{16}\{\varphi(x+\Delta)+\varphi(x-\Delta\}
-\frac{1}{16}\{\varphi(x+3\Delta)+\varphi(x-3\Delta)\}
\]
\noindent der ich folgende Form gebe, die ohne weitere Erklärung verständlich sein
wird:
\clearpage\SourceLeaf{628}\ArithmeticHead{73}{R}{ANHANG.}
\begin{center}
\begin{tabular}{c@{\qquad}l}
zu interpolierende Reihe&\\
$*$&\\
$a$&\\
$*$&$b+c=A$\\
$b$&\\
$*$&$a+d=A\mp B$\\
$c$&\\
$*$&Glied zwischen $b$ und $c=\dfrac{A\pm\frac{1}{8}B}{2}$\\
$d$&\\
$*$&\\
etc.&
\end{tabular}
\end{center}
\noindent\hspace*{\parindent}Hierbei werden nur die vierten Differenzen vernachlässigt, welches hier
völlig erlaubt ist.

Bei diesem Interpolationsgeschäfte müssen dann auch etwaige Fehler,
wenn sie anders mehr als ein Paar Einheiten betragen, sich von selbst zeigen.
Auf diese Art fand ich, dass bei Ihrer Rechnung für $22^0\,12'$ ein Versehen
vorgefallen sein müsse, und indem ich für diesen Winkel die Rechnung selbst
wiederholte, erhielt ich die Logarithmen von $B'''$ an um 21 Unitäten grösser.
Gelegentlich sehen Sie wohl Ihre Rechnung dafür einmal nach. . . .
\begin{center}[4.]\\[3pt]
\textsc{Gauss} an \textsc{Encke}. Göttingen, den 11. März 1851.\end{center}
\noindent\hspace*{\parindent}. . . Der junge \textsc{Dase}, dessen ich in dem andern Brief erwähne, war im
vorigen Herbst hier, hat aber meiner Erwartung keineswegs entsprochen. Seine
Fertigkeit des schnellen Zählens, z. B. von Linsen u. dgl. ist ein sehr merkwürdiges
Talent, aber in keiner seiner das Rechnen betreffenden Leistungen
finde ich etwas wunderbar. Sein Bravourstück, das Multiplizieren zweier vielziffriger
Zahlen \textls[120]{im Kopfe} ist gar nicht wunderbar, weil es in der Tat gar
\textls[120]{kein Rechnen im Kopfe ist}. Er muss notwendig fortwährend die beiden
Faktoren vor sich stehen haben. Allerdings wenn man hört, er multipliziere
eine 100ziffrige Zahl mit einer eben so langen andern \textls[120]{im Kopfe} (in $8^{\mathrm h}\,45'$
in München, wie er jedem erzählt) und meint, er habe sämtliche Operationen
bloss in seiner Phantasie vor sich, ohne die Faktoren länger als ein paar
\par\noindent{\footnotesize X 2 Abh. 6.\hfill 10}
\clearpage\SourceLeaf{629}\ArithmeticHead{74}{L}{PHILIPP MAENNCHEN, GAUSS ALS ZAHLENRECHNER.}
\setcounter{footnote}{0}
\noindent Sekunden gesehen zu haben, so erscheint dies miraculös. Aber die Sache ist
ja gar nicht so. Da die Faktoren ihm fortwährend vor Augen stehen, z. B.
\[
\begin{array}{ccccccccc}
. & . & . & . & . & . & c & b & a\\
. & . & . & . & . & . & \gamma & \beta & \alpha
\end{array}
\]
\noindent so rechnet er $a\alpha$, schreibt die letzte Ziffer hin und addiert die erste mit $a\beta$
und $b\alpha$ zusammen, schreibt von der Summe wieder nur die letzte Ziffer hin
und addiert die erste successive mit $a\gamma,\,b\beta,\,c\alpha$ zusammen u. s. w.\begingroup[\footnotemark]\renewcommand{\thefootnote}{[\arabic{footnote})}\footnotetext{Man erkennt hier, dass \textsc{Gauss} auch die sogenannte symmetrische Multiplikation bekannt und
geläufig war. Den \textls[120]{praktischen} Wert dieses Verfahrens scheint er nicht hoch eingeschätzt zu haben.]}\endgroup. Man
braucht also jedesmal nur sehr wenige Ziffern auf ganz kurze Zeit im Gedächtnis
zu behalten und jedes Kind, welches das Einmaleins kann, lernt es
bald ebenso gut, wenn auch nicht ebenso schnell ausführen. Seine Rechenfertigkeit
beschränkt sich auf die mechanische Ausübung der vier Spezies. Bei
einer (an sich ganz leichten) Aufgabe, wo wiederholt Rechnung mit anderweitiger
Überlegung abwechselt, wird er leicht konfus und gibt wohl ganz
unrichtige Resultate.

Übrigens dies Urteil nur unter uns. Ich fürchte, dass die vielen ungeschickten
Lobhudler in seinem Album ihm nur einen schlechten Dienst geleistet
haben, indem sie ihm eine ganz törichte Meinung von der Bedeutung
seiner Fertigkeiten beigebracht und befestigt haben. . . .
\begin{center}[5.]\\[3pt]
\textsc{Gauss} an \textsc{Olbers}. 13. Februar 1821.\end{center}
\noindent\hspace*{\parindent}. . . Aber auch nachher habe ich mich nur auf die Kometenbeobachtungen
und die Durchgänge von ein paar Sternen beschränkt, indem ich die
zufällige Unterbrechung einer meiner Vorlesungen benutzt habe, wieder eine
theoretische Arbeit vorzunehmen, die ich schon 1818 angefangen, aber bei
meiner zerstückelten Zeit und so mannigfaltigen zum Teil widerwärtigen und
nicht immer die zu \textls[120]{solchen} Arbeiten nötige freie Heiterkeit des Geistes
lassenden Beschäftigungen oft auf lange Zeit wieder weggelegt hatte. Es ist
die neue Begründung der sogenannten Methode der kleinsten Quadrate oder
vielmehr eine ziemlich ausgedehnte allgemeine Untersuchung, wovon diese
\clearpage\SourceLeaf{630}\ArithmeticHead{75}{R}{ANHANG.}
\noindent nur Ein Teil ist. Jetzt ist die erste Hälfte ganz vollendet, die ich in kurzem
der Sozietät zu übergeben denke; die zweite, welche auch bis auf einiges
noch zu überarbeitende fertig ist, wird vermutlich auch noch vor Ostern mitgedruckt
werden können. Sie werden manche artige Sachen darin finden.
Mit Betrübnis fühle ich, wie wenig ich in meiner Lage mit \textls[120]{allen} ihren Missverhältnissen
von dem leisten kann, was ich vielleicht unter glücklicheren Umständen
hätte leisten können, und dass wohl selbst der grössere Teil meiner
früheren Lukubrationen mit mir untergehen wird. — Verzeihen Sie, teuerster
\textsc{Olbers}, den Ausdruck eines Gefühls, welches gerade jetzt beim Empfang
eines mit jugendlichem Feuer geschriebenen Briefes von einem 18jährigen
Florentiner, namens \textsc{Libri}, der mir eine kleine vielversprechende Abhandlung
über höhere Arithmetik zuschickte, wieder recht lebendig bei mir geworden
ist. . . .
\begin{center}\rule{.24\linewidth}{.35pt}\end{center}
\clearpage\endgroup

